\section{Open Questions}\label{sec:oq}
\begin{description}
  \item[Q1.] \textbf{Does the $O(1)$ overhead survive a \emph{chained} deferred-measurement
        construction on $\Omega(\log n)$ intermediate measurements?}  On our two toy circuits
        we observed 0--1 ancilla overhead per protocol, matching the paper's claim.  But
        Fefferman-Lin's argument crucially involves \emph{iterated} purification/amplitude
        amplification with many intermediate-measurement layers.  Empirically the naive
        deferred-measurement rewrite would need one fresh ancilla per measurement layer,
        which is $\Theta(T)$ not $O(1)$.  A concrete follow-up is to build a 10-layer RUS
        pipeline, apply the paper's Sec.~5.2 purification-and-reuse trick (rewind + reuse
        ancilla), and check that (a) the ancilla count stays $\le 2$ regardless of depth
        and (b) statevector-precision equivalence is preserved.
        \textit{Basis:} we tested single-layer circuits only; the O(1) claim is aggregate,
        not per-layer.
  \item[Q2.] \textbf{How does the constant hidden inside the $O(k)$ in Theorem 19's classical
        space-reduction blow up on realistic instances?}  The paper is asymptotic; the
        implied constant on the reduction from Well-Conditioned-Matrix-Inversion to a
        space-$O(k)$ circuit is not tracked.  For $k=6$ (i.e.\ inverting a $64\times 64$
        well-conditioned matrix on 6 logical qubits + ancillas), how many total qubits does
        the naive implementation actually need?  10?  50?  200?  This bears directly on
        whether the paper's construction is a real algorithmic recipe or purely a
        complexity-theoretic scaffold.
        \textit{Basis:} the paper never gives a concrete resource count for any $k$.
  \item[Q3.] \textbf{Does deferred measurement preserve error robustness under a realistic
        noise channel, or does it amplify decoherence?}  Our reproduction uses noise-free
        Aer statevector.  In practice, holding the would-be-measured ancillas to the end
        keeps them entangled with the data during the whole circuit and exposes them to
        depolarizing / amplitude-damping noise for $T$ additional gate cycles.  Repeat
        Experiment A under a depolarizing channel of strength $p$ and compare
        end-of-circuit TV vs.\ mid-measurement-with-classical-control TV for $p \in
        \{10^{-4}, 10^{-3}, 10^{-2}\}$.  Hypothesis: deferred is strictly worse by a factor
        of the ancilla depth.
        \textit{Basis:} Fefferman-Lin do not discuss any noise resilience; the theorem is
        for ideal unitary circuits.  A negative answer here would be a real limitation of
        the ``unitary space is enough'' story in the NISQ era.
  \item[Q4.] \textbf{Is there a matching \emph{lower} bound on the ancilla overhead --- i.e.\
        is $O(1)$ tight?}  The paper proves $\le O(1)$.  It is unclear whether any circuit
        exists that provably requires $\Omega(1)$ ancillas (as opposed to $0$).  Our
        teleportation example achieves $0$ and RUS achieves $1$; is there an intermediate-
        measurement circuit --- e.g.\ a magic-state distillation subroutine --- for which
        no purification/reuse scheme can go below $\ge 2$ ancillas?  A separation result
        here would sharpen the constant and give the community a canonical hard example.
        \textit{Basis:} we observed different overheads for two structurally different
        circuits (0 vs.\ 1); the question is whether the achievable minimum depends on the
        circuit structure or is always $\le 1$.
  \item[Q5.] \textbf{Does Fefferman-Lin's construction extend to \emph{measurement-and-adaptive-
        gate-choice} circuits (dynamic circuits) rather than only measurement-and-Pauli-correction
        circuits?}  Our reproductions used only Pauli corrections ($X$, $Z$).  Real dynamic-circuit
        primitives on IBM and Quantinuum hardware use classically-computed \emph{arbitrary}
        single-qubit rotations conditional on prior outcomes (e.g.\ classically-computed $R_z(\theta)$
        with $\theta$ depending on prior measurement bits).  Is the same $O(1)$ ancilla
        upper bound achievable in this richer model, or does the classical-computation-of-
        the-angle step break the trick?  Repeat the experiment with an $R_z(2\pi \cdot 0.\!m_0m_1)$
        (i.e.\ quantum Fourier transform-style feedback) correction and check whether a
        finite-ancilla coherent implementation exists.
        \textit{Basis:} the paper's construction is stated for the classical
        Pauli-correction model.  Modern dynamic circuits routinely use richer classical
        computation, and the ``$O(1)$ overhead'' story may not transfer.
\end{description}
